% QMB_n_ai.tex — Appendix: Note on AI Assistance
\chapter*{Note on the Origin}
\addcontentsline{toc}{chapter}{Note on the Origin}

This book grew out of the FFGFT corpus, which has comprised over 365
individual documents since 2025. The texts of the individual chapters were
conceived and substantively overseen by Johann Pascher; Claude
(Anthropic, Sonnet~4.6) assisted with the formulation, with the merging
of the source sections, and with the compilation of the \LaTeX template.

\medskip\noindent\textbf{What AI assistance meant in this project:}
The source documents (Doc.~022 to Doc.~365) contain the actual
derivations, test scripts, and epistemic markers. The AI read these
documents, reformulated the chapters in its own language --- without
verbatim copying, without external content --- and checked the consistency
of the status markers. All substantive decisions, boundary definitions, and
status markers are Johann's work; the AI was a formulation tool,
not a substantive author.

\medskip\noindent\textbf{What AI assistance did not mean:}
None of the physical statements originate from the AI's training corpus.
All numbers, derivations, and cross-references trace back to the
aforementioned source documents, which are publicly viewable in the repository.
Questions as to whether a result is \enquote{new} or \enquote{correct}
were answered by the author himself, not by the AI.

\medskip\noindent\textbf{Verifiability:} Each chapter bears at the end
a \texttt{\\quelle} reference to the underlying
corpus documents. Each substantive statement bears a status marker
([K], [B], [S], [E], [X]). Anyone who wishes to question a statement
will find the complete derivation path with
test script in the cited document.